\chapter{Operator K-theory and the analytic side}
\label{ch:operator-k-theory}

Unless stated otherwise, $A$ is a unital complex $C^*$-algebra.
The present Lean milestone constructs projections, fixed-size equivalence
classes, matrix block sums, and norm homotopies of matrix unitaries.
The stable quotients and the groups $K_0$ and $K_1$ remain separate targets.

\section{Projections and Murray--von Neumann equivalence}

\begin{definition}\label{def:operator-projection}
\lean{BC4lean.OperatorKTheory.Projection,BC4lean.OperatorKTheory.Projection.zero,BC4lean.OperatorKTheory.Projection.one,BC4lean.OperatorKTheory.Projection.complement}
\leanok
A projection in $A$ is an element $p$ satisfying $p^*=p$ and $p^2=p$.
Write $\operatorname{Proj}(A)$ for these projections. The elements $0$, $1$,
and $1-p$ are projections.
\end{definition}

\begin{proposition}\label{prop:partial-isometry-support}
\uses{def:operator-projection}
\lean{BC4lean.OperatorKTheory.Projection.right_support}
If $p\in\operatorname{Proj}(A)$ and $v^*v=p$, then $vp=v$.
\end{proposition}
\begin{proof}
\leanok
For $x=vp-v$, expansion using $p^*=p$, $p^2=p$, and $v^*v=p$
gives $x^*x=0$. The $C^*$-identity implies $x=0$.
\end{proof}

\begin{definition}\label{def:projection-equivalence}
\uses{def:operator-projection}
\lean{BC4lean.OperatorKTheory.Projection.Equivalent}
\leanok
Projections $p,q\in A$ are Murray--von Neumann equivalent, written $p\sim q$,
if some $v\in A$ satisfies $v^*v=p$ and $vv^*=q$.
\end{definition}

\begin{proposition}\label{prop:projection-equivalence}
\uses{def:projection-equivalence,prop:partial-isometry-support}
\lean{BC4lean.OperatorKTheory.Projection.equivalent_refl,BC4lean.OperatorKTheory.Projection.equivalent_symm,BC4lean.OperatorKTheory.Projection.equivalent_trans,BC4lean.OperatorKTheory.Projection.equivalenceSetoid,BC4lean.OperatorKTheory.Projection.Classes}
The relation $\sim$ is an equivalence relation, so its quotient is defined.
\end{proposition}
\begin{proof}
\leanok
Reflexivity uses $v=p$ and symmetry uses $v^*$. If $v$ implements $p\sim q$
and $w$ implements $q\sim r$, then $wv$ implements $p\sim r$.
The support identities $vp=v$, $qv=v$, and $wq=w$ give the two required equations.
These quotients concern one algebra at a time; they are not yet stable classes.
\end{proof}

\begin{proposition}\label{prop:projection-map}
\uses{prop:projection-equivalence}
\lean{BC4lean.OperatorKTheory.Projection.map,BC4lean.OperatorKTheory.Projection.equivalent_map,BC4lean.OperatorKTheory.Projection.mapClasses}
A unital complex star homomorphism $\phi:A\to B$ sends projections to
projections and preserves Murray--von Neumann equivalence. It therefore
induces a map on the preceding quotients.
\end{proposition}
\begin{proof}
\leanok
Apply $\phi$ to the projection equations and, for equivalence, to the
witness $v$ and its two defining equations.
\end{proof}

\section{Matrix stabilization and the construction of $K_0$}

\begin{definition}\label{def:cstar-matrix-projection}
\uses{def:operator-projection}
\lean{BC4lean.OperatorKTheory.MatrixProjection,BC4lean.OperatorKTheory.matrixBlockSum,BC4lean.OperatorKTheory.projectionBlockSum,BC4lean.OperatorKTheory.stabilizeProjection}
\leanok
Use the $C^*$-norm on $M_n(A)$, and let $P_n(A)=\operatorname{Proj}(M_n(A))$.
For $p\in P_n(A)$ and $q\in P_m(A)$, their block sum is
\[
 p\oplus q=\begin{pmatrix}p&0\\0&q\end{pmatrix}\in P_{n+m}(A).
\]
Projection stabilization is $p\mapsto p\oplus 0_m$.
The Lean implementation initially uses arbitrary finite index types and their
disjoint sums. Mathlib's matrix norm requires a compatible star order on $A$;
the matrix declarations retain this explicit implementation assumption.
\end{definition}

\begin{proposition}\label{prop:projection-block-sum}
\uses{def:cstar-matrix-projection,prop:projection-equivalence}
\lean{BC4lean.OperatorKTheory.equivalent_blockSum,BC4lean.OperatorKTheory.equivalent_stabilize,BC4lean.OperatorKTheory.blockSumClasses}
If $p\sim p'$ and $q\sim q'$, then $p\oplus q\sim p'\oplus q'$.
Thus block sum and stabilization descend to equivalence classes at the
indicated matrix sizes.
\end{proposition}
\begin{proof}
\leanok
If $v$ and $w$ are the respective witnesses, use $v\oplus w$.
Block multiplication and adjoints give both defining equations.
\end{proof}

\begin{definition}\label{def:stable-projection-monoid}
\uses{prop:projection-block-sum}
\notready
On the disjoint union of the $P_n(A)$, identify $p\in P_n(A)$ and
$q\in P_m(A)$ when, for some $N\geq n,m$,
\[
 p\oplus 0_{N-n}\sim q\oplus 0_{N-m}\quad\text{in }M_N(A).
\]
Construct this stable quotient $V(A)$ and prove that block sum gives a
commutative monoid, independently of representatives and finite-index
identifications. Its zero is the class of a zero projection.
\end{definition}

\begin{definition}\label{def:operator-k0}
\uses{def:stable-projection-monoid}
\notready
Define $K_0(A)$ as the Grothendieck group of $V(A)$, with its canonical map
$V(A)\to K_0(A)$. Prove its universal property for monoid maps from $V(A)$
to abelian groups. Do not assume that the canonical map is injective:
cancellation in $V(A)$ is not a standing hypothesis.
\end{definition}

\section{Unitaries, norm homotopy, and $K_1$}

\begin{definition}\label{def:matrix-unitary}
\lean{BC4lean.OperatorKTheory.MatrixUnitary,BC4lean.OperatorKTheory.unitaryBlockSum,BC4lean.OperatorKTheory.stabilizeUnitary}
\leanok
Let $U_n(A)=\{u\in M_n(A):u^*u=uu^*=1\}$, with the subspace topology
from the matrix $C^*$-norm. Block sums of unitaries are unitary.
Unitary stabilization is $u\mapsto u\oplus 1_m$.
\end{definition}

\begin{definition}\label{def:unitary-norm-homotopy}
\uses{def:matrix-unitary}
\lean{BC4lean.OperatorKTheory.UnitaryHomotopic,BC4lean.OperatorKTheory.UnitaryHomotopyClasses}
\leanok
Two elements of $U_n(A)$ are norm homotopic if they are endpoints of a
continuous path $[0,1]\to U_n(A)$. Form the quotient by this equivalence
relation at each fixed size.
\end{definition}

\begin{proposition}\label{prop:unitary-homotopy-operations}
\uses{def:unitary-norm-homotopy}
\lean{BC4lean.OperatorKTheory.unitaryHomotopic_refl,BC4lean.OperatorKTheory.unitaryHomotopic_symm,BC4lean.OperatorKTheory.unitaryHomotopic_trans,BC4lean.OperatorKTheory.unitaryHomotopic_mul,BC4lean.OperatorKTheory.unitaryHomotopic_inv}
Norm homotopy is an equivalence relation and is preserved by multiplication
of pairs of unitaries and by inversion.
\end{proposition}
\begin{proof}
\leanok
Use constant paths, reversal, concatenation, pointwise multiplication,
and pointwise inversion in the topological unitary group.
\end{proof}

\begin{proposition}\label{prop:unitary-homotopy-stabilization}
\uses{prop:unitary-homotopy-operations,def:matrix-unitary}
\lean{BC4lean.OperatorKTheory.continuous_matrixBlockSum,BC4lean.OperatorKTheory.continuous_unitaryBlockSum,BC4lean.OperatorKTheory.unitaryHomotopic_blockSum,BC4lean.OperatorKTheory.unitaryHomotopic_stabilize,BC4lean.OperatorKTheory.stabilizeUnitaryClasses}
Block sum is continuous and preserves norm homotopies. Consequently,
stabilization induces maps on the fixed-size norm-path components.
\end{proposition}
\begin{proof}
\leanok
The finite matrix topology agrees with the entrywise product topology,
in which block sum is continuous. Apply block sum to a pair of paths;
for stabilization, use a constant identity path in the second factor.
\end{proof}

\begin{definition}\label{def:operator-k1}
\uses{prop:unitary-homotopy-stabilization}
\notready
Take the stable quotient of matrix unitaries by norm homotopy after adjoining
identity blocks. Define $K_1(A)$ to be this quotient and prove that block sum
gives an abelian group. This requires the stabilization and reindexing
coherence, block-switching homotopies, and an inverse represented by $u^*$.
The fixed-size quotients above do not yet establish these assertions.
\end{definition}

\section{Nonunital algebras and structural properties}

\begin{definition}\label{def:operator-k-theory}
\uses{def:operator-k0,def:operator-k1}
\notready
For a possibly nonunital complex $C^*$-algebra $A$, use the forced unitization
$A^+$ and its scalar quotient $\epsilon:A^+\to\mathbb C$ to define
\[
 K_i(A)=\ker\bigl(K_i(\epsilon):K_i(A^+)\to K_i(\mathbb C)\bigr),
 \qquad i=0,1.
\]
Construct the induced homomorphisms needed here and prove agreement with
the preceding definitions when $A$ is unital.
\end{definition}

\begin{proposition}\label{prop:operator-k-functoriality}
\uses{def:operator-k-theory,prop:projection-map}
\notready
Complex star homomorphisms, including nonunital ones, induce homomorphisms
on $K_i$ that preserve identities and composition.
\end{proposition}

\begin{proposition}\label{prop:operator-k-homotopy-invariance}
\uses{prop:operator-k-functoriality}
\notready
Point-norm homotopic star homomorphisms induce the same maps on $K_i$.
\end{proposition}

The six-term exact sequence and Bott periodicity require separate future
nodes and proofs. They are not consequences of the formalized foundations
merely by being part of the usual theory.

\section{The analytic target}

\begin{definition}\label{def:analytic-side}
\uses{def:reduced-group-cstar,def:operator-k-theory}
\notready
For a countable discrete group $\Gamma$, the analytic target of the
coefficient-free reduced Baum--Connes assembly map is
\[
 K_i(C_r^*(\Gamma)),\qquad i=0,1.
\]
The reduced algebra has already been constructed. Instantiating the
$K$-theory functors and connecting them to the assembly map remain to be done.
\end{definition}
